\documentclass[11pt]{article}

\usepackage{lmodern}   % Type 1 T1 fonts; without this T1 falls back to EC bitmaps
\usepackage[T1]{fontenc}
\usepackage[utf8]{inputenc}
\usepackage[margin=1in]{geometry}
\usepackage{amsmath}
\usepackage{graphicx}
\usepackage{float}
\usepackage{booktabs}
\usepackage{longtable}
\usepackage{enumitem}
\usepackage{hyperref}
\usepackage{xcolor}
\usepackage{listings}
\usepackage{textcomp}
\usepackage{upquote}  % straight quotes in listings/verbatim (needs T1 + textcomp)

\hypersetup{
  colorlinks=true,
  linkcolor=blue,
  urlcolor=blue,
  pdftitle={CTA Live Wall Map Build Guide},
  pdfauthor={DIY Transit Map},
}

% ---------- Code listing styles ----------
\lstdefinestyle{arduino}{
  language=C++,
  basicstyle=\ttfamily\small,
  keywordstyle=\color{blue}\bfseries,
  commentstyle=\color{gray},
  stringstyle=\color{orange},
  breaklines=true,
  showstringspaces=false,
  upquote=true,
  frame=single,
  numbers=left,
  numberstyle=\tiny,
  captionpos=b,
  morekeywords={String, boolean, byte, unsigned, long, uint8_t, uint16_t, uint32_t, CRGB, JsonArray, JsonObject, DynamicJsonDocument, DeserializationError, HTTPClient, WiFiClientSecure}
}

\lstdefinestyle{python}{
  language=Python,
  basicstyle=\ttfamily\small,
  keywordstyle=\color{blue}\bfseries,
  commentstyle=\color{gray},
  stringstyle=\color{orange},
  breaklines=true,
  showstringspaces=false,
  upquote=true,
  frame=single,
  numbers=left,
  numberstyle=\tiny,
  captionpos=b
}

\lstdefinestyle{text}{
  basicstyle=\ttfamily\small,
  breaklines=true,
  showstringspaces=false,
  upquote=true,
  frame=single,
  captionpos=b
}

\title{\textbf{Build a Live Chicago CTA Wall Map}\\\large A Step-by-Step Guide for a Cheap ESP32 + WS2812B DIY Project}
\author{DIY Transit Electronics}
\date{\today}

\begin{document}

\maketitle

\tableofcontents
\newpage

\section{Introduction}

This guide describes how to build a framed wall map of the Chicago Transit Authority (CTA) ``L'' rail system where a small RGB LED at each station lights up when a train is approaching or has just departed that station. The data comes live from the official CTA Train Tracker API.

The build is intended to be as cheap as possible while still producing a reliable, attractive object. It uses:

\begin{itemize}
    \item An ESP32 development board as the WiFi microcontroller.
    \item WS2812B addressable RGB LEDs, one per station.
    \item A printed or laser-cut schematic map mounted in a shadow box.
    \item A 5V power supply.
    \item Free software: Arduino IDE, Python, Inkscape or similar.
\end{itemize}

You should have basic soldering skills and be comfortable uploading code to an ESP32. This project is for \textbf{personal use}. The CTA map, logo, and data have copyright/trademark restrictions; do not sell the finished item without addressing licensing and API terms.

\section{Bill of Materials}

The total can be around \$70--\$140 depending on what you already own. The table below gives the recommended cheap parts.

% Column widths sum to 14.2cm; with 10 x \tabcolsep (2.1cm) this fits the
% 16.5cm textwidth left by the 1in margins.
\begin{longtable}{p{3.2cm} p{4.6cm} p{1.4cm} p{2.0cm} p{3.0cm}}
\toprule
\textbf{Item} & \textbf{Specification} & \textbf{Qty} & \textbf{Approx. Cost} & \textbf{Source} \\
\midrule
ESP32 Dev Board & WROOM-32, USB-C or micro-USB, 30-pin or 38-pin & 1 & \$4--\$8 & AliExpress, Amazon, eBay \\
WS2812B LED strip & 30 LED/m, 5V, IP30 non-waterproof, 5m roll (150 LEDs) & 1 & \$12--\$20 & AliExpress, Amazon \\
Power supply & 5V DC, 10A, barrel jack or screw terminals & 1 & \$10--\$15 & Amazon, AliExpress \\
Level shifter & 74AHCT125 or 74HCT245 DIP & 1 & \$0.50--\$2 & eBay, local electronics \\
330 $\Omega$ resistor & 1/4 W, metal film or carbon & 1 & pennies & local or online \\
1000 $\mu$F capacitor & 10V or higher electrolytic & 1 & \$0.50 & local or online \\
Hook-up wire & 22 AWG solid core, multiple colors & 5--10 m & \$5 & local or online \\
Heat shrink tubing & Assorted small diameters & 1 pack & \$3 & local or online \\
Foam board or plywood & 24$\times$36 inch or larger & 1 & \$3--\$10 & dollar store, art store, hardware store \\
Printed map & 24$\times$36 inch color print of a simplified CTA L map & 1 & \$10--\$25 & local print shop, FedEx Office \\
Diffuser material & tracing paper, vellum, or thin white acrylic & 1 sheet & \$2--\$8 & art store or online \\
Frame/shadow box & 24$\times$36 inch or custom & 1 & \$15--\$40 & IKEA, thrift store, craft store \\
Solder, flux, glue, tape & Misc consumables & --- & \$10--\$15 & local or online \\
\bottomrule
\caption{Cheap BOM for the CTA wall map. Prices are approximate and vary by region.}
\label{tab:bom}
\end{longtable}

\subsection*{Tools You Will Need}

\begin{itemize}
    \item Soldering iron with fine tip
    \item Solder wire and flux
    \item Wire strippers
    \item Multimeter
    \item Hot glue gun or double-sided tape
    \item Computer with Arduino IDE
    \item USB cable for ESP32
    \item Hobby knife or scissors
    \item Ruler or straightedge
\end{itemize}

\section{Step 1: Get the Data Feeds}

\subsection{1.1 CTA Train Tracker API Key}

The live train positions come from CTA's official API. You need a free developer key.

\begin{enumerate}
    \item Go to \url{https://www.transitchicago.com/developers/traintracker/}.
    \item Register for a developer account if you do not have one.
    \item Request access to the \textbf{Train Tracker API}.
    \item You will receive an API key, usually a string of letters and numbers.
\end{enumerate}

The key is free for non-commercial use. Read the terms of service. Do not share your key publicly.

\subsection{1.2 CTA GTFS Static Data}

You also need the static schedule and station data, provided in GTFS format.

\begin{itemize}
    \item Download from: \url{https://www.transitchicago.com/downloads/sch_data/google_transit.zip}
    \item Extract the ZIP file somewhere on your computer.
\end{itemize}

The GTFS file contains \texttt{stops.txt}, \texttt{routes.txt}, \texttt{trips.txt}, and \texttt{stop\_times.txt}, which we will use to build a station list.

\section{Step 2: Generate the Station List}

We need a list of every rail station with its CTA \texttt{stop\_id}. The script below downloads the GTFS feed, filters for rail routes, and writes a CSV template.

Create a file named \texttt{prepare\_gtfs.py}:

\begin{lstlisting}[style=python,caption={Python script to extract CTA rail stations from GTFS.},label={lst:prepare_gtfs}]
import urllib.request
import zipfile
import csv
import io
import os
from collections import defaultdict

GTFS_URL = "https://www.transitchicago.com/downloads/sch_data/google_transit.zip"
OUT_DIR = "cta_gtfs"
OUT_CSV = "stations_template.csv"

def download_gtfs():
    print("Downloading GTFS feed...")
    data = urllib.request.urlopen(GTFS_URL).read()
    with zipfile.ZipFile(io.BytesIO(data)) as z:
        z.extractall(OUT_DIR)
    print(f"Extracted GTFS to {OUT_DIR}")

def read_csv(filename):
    with open(filename, newline='', encoding='utf-8-sig') as f:
        return list(csv.DictReader(f))

def main():
    if not os.path.exists(os.path.join(OUT_DIR, "stops.txt")):
        download_gtfs()

    routes = read_csv(os.path.join(OUT_DIR, "routes.txt"))
    trips = read_csv(os.path.join(OUT_DIR, "trips.txt"))
    stop_times = read_csv(os.path.join(OUT_DIR, "stop_times.txt"))
    stops = read_csv(os.path.join(OUT_DIR, "stops.txt"))

    # Filter for rail routes. In CTA GTFS, route_type == 1 means subway/metro.
    rail_routes = [r for r in routes if r.get("route_type") == "1"]
    if not rail_routes:
        # Fallback: use known route short names.
        known = {"Red", "Blue", "Brn", "G", "Org", "Pink", "P", "Y"}
        rail_routes = [r for r in routes if r.get("route_short_name") in known]

    route_ids = {r["route_id"] for r in rail_routes}
    route_short_names = {r["route_id"]: r["route_short_name"] for r in rail_routes}

    # Map trip_id to route_id for rail trips only.
    trip_to_route = {}
    for t in trips:
        if t["route_id"] in route_ids:
            trip_to_route[t["trip_id"]] = t["route_id"]

    stop_to_routes = defaultdict(set)
    stop_id_to_name = {s["stop_id"]: s["stop_name"] for s in stops}

    # Iterate through stop_times. This may take a minute because the file is large.
    print("Scanning stop_times.txt...")
    for st in stop_times:
        route_id = trip_to_route.get(st["trip_id"])
        if route_id is not None:
            stop_id = st["stop_id"]
            stop_to_routes[stop_id].add(route_id)

    print(f"Found {len(stop_to_routes)} unique rail station IDs.")

    # Write CSV template. You will fill in led_index manually later.
    with open(OUT_CSV, "w", newline='', encoding="utf-8") as f:
        writer = csv.writer(f)
        writer.writerow(["line", "stop_id", "stop_name", "led_index"])
        for stop_id, routes_for_stop in sorted(stop_to_routes.items()):
            lines = [route_short_names[rid] for rid in routes_for_stop]
            writer.writerow([",".join(lines), stop_id, stop_id_to_name.get(stop_id, ""), ""])

    print(f"Wrote {OUT_CSV}")
    print("Open the CSV and fill in the led_index column.")

if __name__ == "__main__":
    main()
\end{lstlisting}

Run the script:

\begin{verbatim}
python prepare_gtfs.py
\end{verbatim}

This creates \texttt{stations\_template.csv} with columns:

\begin{verbatim}
line,stop_id,stop_name,led_index
Red,41090,Jarvis,
Red,41100,Morse,
...
\end{verbatim}

\section{Step 3: Design the Physical Map and Assign LED Indices}

You need a physical map with one LED per station. A geographically accurate map will be too crowded downtown, so a schematic transit map is better.

\subsection{3.1 Create or Obtain a Map}

For a personal project you can trace a simplified CTA map, draw one in Inkscape, or use a map image for personal reference. Keep these points in mind:

\begin{itemize}
    \item Leave enough space between Loop stations.
    \item Use a dark background; black or dark charcoal makes LED light pop.
    \item Mark each station with a small circle or dot.
    \item Do not use the official CTA logo if you might later sell the item.
\end{itemize}

A good size is 24$\times$36 inches. You can print the map at a local print shop.

\subsection{3.2 Mount the Map and Mark Holes}

Attach the printed map to a rigid board such as foam board or thin plywood. Poke or drill a small hole at each station location. The holes should be slightly smaller than the LED package diameter, typically 8--10~mm for WS2812B 5050 LEDs.

\subsection{3.3 Number Every Station}

Starting at one end of the map, assign each station a unique integer LED index beginning at 0. Write the number next to each hole on the back of the board. This number is the \texttt{led\_index} in the CSV.

Be systematic. A suggested order:

\begin{enumerate}
    \item Red Line from Howard to 95th/Dan Ryan
    \item Blue Line from O'Hare to Forest Park
    \item Brown Line
    \item Green Line
    \item Orange Line
    \item Pink Line
    \item Purple Line
    \item Yellow Line
    \item Shared Loop stations only once, with one LED per physical station.
\end{enumerate}

Make sure no two stations share the same index. The total should be around 145--146 unique stations.

\subsection{3.4 Fill in the CSV}

Open \texttt{stations\_template.csv} in a spreadsheet program. For each row, enter the \texttt{led\_index} that you assigned to that station. Save the file.

\section{Step 4: Generate the Arduino Mapping Header}

The ESP32 firmware needs a mapping from CTA \texttt{stop\_id} to LED index. The script below reads your filled CSV and emits a C++ header file.

Create \texttt{make\_mapping.py}:

\begin{lstlisting}[style=python,caption={Python script to generate mapping.h from the filled CSV.},label={lst:make_mapping}]
import csv
import sys

def main(csv_file, output_header):
    rows = []
    max_led = -1

    with open(csv_file, newline='', encoding='utf-8') as f:
        reader = csv.DictReader(f)
        for row in reader:
            led_str = row.get("led_index", "").strip()
            stop_id = row.get("stop_id", "").strip()
            if led_str == "" or stop_id == "":
                continue
            led_index = int(led_str)
            rows.append((led_index, stop_id))
            max_led = max(max_led, led_index)

    if not rows:
        print("Error: no rows with led_index filled in.")
        sys.exit(1)

    rows.sort()

    with open(output_header, "w", encoding="utf-8") as f:
        f.write("#ifndef MAPPING_H\n")
        f.write("#define MAPPING_H\n\n")
        f.write(f"#define NUM_LEDS {max_led + 1}\n")
        f.write(f"const int stationCount = {len(rows)};\n\n")
        f.write("struct StationMap {\n")
        f.write("  const char* stopId;\n")
        f.write("  int ledIndex;\n")
        f.write("};\n\n")
        f.write("const StationMap stationMap[] = {\n")
        for led_index, stop_id in rows:
            f.write(f'  {{"{stop_id}", {led_index}}},\n')
        f.write("};\n\n")
        f.write("#endif\n")

    print(f"Wrote {output_header} with {len(rows)} stations, NUM_LEDS={max_led + 1}")

if __name__ == "__main__":
    if len(sys.argv) != 3:
        print("Usage: python make_mapping.py stations_template.csv mapping.h")
        sys.exit(1)
    main(sys.argv[1], sys.argv[2])
\end{lstlisting}

Run it:

\begin{verbatim}
python make_mapping.py stations_template.csv mapping.h
\end{verbatim}

This creates \texttt{mapping.h}. An example of the first few lines:

\begin{lstlisting}[style=arduino,caption={Example generated mapping.h (fragment).},label={lst:mapping_example}]
#ifndef MAPPING_H
#define MAPPING_H

#define NUM_LEDS 145
const int stationCount = 145;

struct StationMap {
  const char* stopId;
  int ledIndex;
};

const StationMap stationMap[] = {
  {"41090", 0},
  {"41100", 1},
  // ... many more ...
};

#endif
\end{lstlisting}

\section{Step 5: Prototype with a Few LEDs}

Before building the full map, test the electronics and firmware with a small strip of 5--10 LEDs.

\begin{enumerate}
    \item Cut a short segment of WS2812B strip containing 5--10 LEDs.
    \item Solder wires to the 5V, GND, and DIN pads on the input end of the strip.
    \item Connect the ESP32 to the strip according to the wiring diagram in Section~\ref{sec:wiring}.
    \item Install the required Arduino libraries.
    \item Use the firmware in Section~\ref{sec:firmware}, but temporarily set \texttt{NUM\_LEDS} to a small number for the test.
    \item Upload the code and verify that the LEDs light up when you call \texttt{FastLED.show()}.
    \item Do not proceed until the prototype works reliably.
\end{enumerate}

\section{Step 6: Wire and Mount the Full LED Array}

\subsection{6.1 Cut the LED Strip}

You need one WS2812B LED per station. Cut the 5m strip into individual LED segments at the copper pads. The data direction is marked by arrows or \texttt{DIN} and \texttt{DOUT} labels.

\begin{itemize}
    \item The input side has \texttt{DIN}, 5V, GND.
    \item The output side has \texttt{DOUT}, 5V, GND.
    \item You only use the \texttt{DIN} of the first LED and \texttt{DOUT} of each LED to the next \texttt{DIN}.
\end{itemize}

\subsection{6.2 Solder Extension Wires Between LEDs}

Because the LEDs must be placed at irregular station positions, you cannot use the strip as a continuous line. For each LED:

\begin{itemize}
    \item Solder three wires to the LED's 5V, GND, and \texttt{DOUT} pads.
    \item Route those wires to the next LED's 5V, GND, and \texttt{DIN} pads.
    \item Use 22 AWG wire. Data wire can be thinner; power wires should be at least 22 AWG.
    \item Keep data wires reasonably short, ideally under 30~cm between LEDs. Use a 330~$\Omega$ resistor in series with the first data input to reduce ringing.
    \item Use heat shrink tubing to cover solder joints.
\end{itemize}

This is tedious but cheap.

\subsection{6.3 Mount LEDs Behind the Map}

Position each LED directly behind its station hole. Hot glue works well. Leave a gap of about 1--2~cm between the LED and the map surface so the light diffuses. You can use a spacer layer such as foam board strips or 3D-printed standoffs.

\subsection{6.4 Add Power Injection}

With 145 WS2812B LEDs, voltage drop along the power wires will cause color and brightness problems. Add power injection every 30--40 LEDs.

\begin{itemize}
    \item Connect the 5V and GND from the power supply to the beginning of the LED chain.
    \item Run additional 5V and GND wires to points after LEDs 30, 60, 90, and 120.
    \item Do not connect data at these injection points; only power and ground.
\end{itemize}

\section{Step 7: Connect the ESP32 and Power Supply}
\label{sec:wiring}

The ESP32 GPIO output is 3.3~V, but WS2812B data inputs officially expect 5~V logic. A 74AHCT125 level shifter converts 3.3~V to 5~V and prevents many flickering problems.

\begin{figure}[H]
\centering
\begin{lstlisting}[style=text,caption={Wiring diagram with 74AHCT125 level shifter.},label={lst:wiring_diagram}]
5V Power Supply (+) ----[1000uF cap]----> LED strip 5V bus
5V Power Supply GND (-) ------------------> LED strip GND bus

ESP32 3.3V -------------------------> 74AHCT125 VCC (pin 14)
ESP32 GND --------------------------> 74AHCT125 GND (pin 7)
ESP32 GPIO12 --> 330 ohm --> 74AHCT125 1A (pin 2)

74AHCT125 1Y (pin 3) ---------------> First LED DIN
74AHCT125 1OE (pin 1) --------------> GND

ESP32 GND --------------------------> LED strip GND bus
\end{lstlisting}
\end{figure}

If you skip the level shifter, connect the ESP32 GPIO through a 330~$\Omega$ resistor directly to the first LED DIN. This often works for short chains, but if the first LED flickers or shows random colors, use the level shifter.

\subsection{Power Supply Notes}

\begin{itemize}
    \item Use a 5V DC supply rated at least 4A for low brightness operation.
    \item A 10A supply is safer for full brightness and future expansion.
    \item Connect the 1000~$\mu$F capacitor across the 5V and GND terminals close to the LED strip. The capacitor side marked with a stripe or minus sign goes to GND.
    \item Do not power the LED strip from the ESP32's 5V pin.
\end{itemize}

\section{Step 8: Install and Configure the Firmware}
\label{sec:firmware}

\subsection{8.1 Install Arduino IDE and Libraries}

\begin{enumerate}
    \item Download and install the Arduino IDE from \url{https://www.arduino.cc/en/software}.
    \item Add ESP32 board support. In Preferences, add the boards manager URL:
    \begin{quote}
    \small\url{https://raw.githubusercontent.com/espressif/arduino-esp32/gh-pages/package_esp32_index.json}
    \end{quote}
    Then install \texttt{esp32} by Espressif Systems.
    \item Install the following libraries from the Library Manager:
    \begin{itemize}
        \item \texttt{FastLED}
        \item \texttt{ArduinoJson}
    \end{itemize}
\end{enumerate}

\subsection{8.2 Create the Sketch}

Create a new Arduino sketch named \texttt{cta\_wall\_map}. Copy the code below into the sketch. Place the generated \texttt{mapping.h} file in the same sketch folder.

\begin{lstlisting}[style=arduino,caption={Complete ESP32 Arduino sketch for the CTA wall map.},label={lst:arduino}]
/*
  CTA Wall Map - ESP32 + WS2812B
  Fetches CTA Train Tracker positions and lights station LEDs.
  Requires mapping.h generated by make_mapping.py
*/

#include <WiFi.h>
#include <WiFiClientSecure.h>
#include <HTTPClient.h>
#include <FastLED.h>
#include <ArduinoJson.h>
#include "mapping.h"

// ---------- User configuration ----------
const char* ssid = "YOUR_WIFI_SSID";
const char* password = "YOUR_WIFI_PASSWORD";
const char* ctaApiKey = "YOUR_CTA_API_KEY";

#define LED_PIN 12
#define BRIGHTNESS 64        // 0-255, lower is safer/cheaper
const char* routeCodes = "red,blue,brn,g,org,pink,p,y";
// ----------------------------------------

CRGB leds[NUM_LEDS];
String apiUrl = "https://lapi.transitchicago.com/api/1.0/ttpositions.aspx?key="
                + String(ctaApiKey) + "&rt=" + String(routeCodes) + "&outputType=json";

WiFiClientSecure client;
HTTPClient http;

int stationToLed(const char* stopId) {
  for (int i = 0; i < stationCount; i++) {
    if (strcmp(stationMap[i].stopId, stopId) == 0) {
      return stationMap[i].ledIndex;
    }
  }
  return -1;
}

CRGB lineColor(const char* line) {
  if (strcmp(line, "red") == 0) return CRGB(198, 12, 48);
  if (strcmp(line, "blue") == 0) return CRGB(0, 161, 222);
  if (strcmp(line, "brn") == 0) return CRGB(98, 54, 27);
  if (strcmp(line, "g") == 0) return CRGB(0, 155, 58);
  if (strcmp(line, "org") == 0) return CRGB(249, 70, 28);
  if (strcmp(line, "pink") == 0) return CRGB(226, 126, 166);
  if (strcmp(line, "p") == 0) return CRGB(82, 35, 152);
  if (strcmp(line, "y") == 0) return CRGB(249, 227, 0);
  return CRGB(255, 255, 255);
}

void connectWiFi() {
  Serial.print("Connecting to WiFi");
  WiFi.mode(WIFI_STA);
  WiFi.setAutoReconnect(true);
  WiFi.begin(ssid, password);
  unsigned long start = millis();
  while (WiFi.status() != WL_CONNECTED && millis() - start < 20000) {
    delay(500);
    Serial.print(".");
  }
  if (WiFi.status() == WL_CONNECTED) {
    Serial.println("\nWiFi connected. IP: " + WiFi.localIP().toString());
  } else {
    Serial.println("\nWiFi connection failed. Will retry in loop.");
  }
}

void fetchAndUpdate() {
  if (WiFi.status() != WL_CONNECTED) {
    connectWiFi();
    if (WiFi.status() != WL_CONNECTED) return;
  }

  client.setInsecure();  // Skip certificate validation; personal prototype only.
  http.begin(client, apiUrl);
  int httpCode = http.GET();

  if (httpCode == HTTP_CODE_OK) {
    // CTA response can be ~50-100KB. 128KB JSON buffer is safe on ESP32.
    DynamicJsonDocument doc(128 * 1024);
    DeserializationError err = deserializeJson(doc, http.getStream());

    if (!err) {
      FastLED.clear();
      JsonArray routeArray = doc["ctatt"]["route"].as<JsonArray>();

      for (JsonObject route : routeArray) {
        const char* line = route["@name"];
        JsonArray trains = route["train"].as<JsonArray>();

        for (JsonObject train : trains) {
          const char* nextStaId = train["nextStaId"];
          int ledIndex = stationToLed(nextStaId);
          if (ledIndex >= 0 && ledIndex < NUM_LEDS) {
            // Overwrites if multiple trains at same station.
            // Optional: add occupancy blending or cycling for shared stations.
            leds[ledIndex] = lineColor(line);
          }
        }
      }

      FastLED.show();
      Serial.println("Updated LEDs.");
    } else {
      Serial.print("JSON parse error: ");
      Serial.println(err.c_str());
    }
  } else {
    Serial.printf("HTTP error: %d\n", httpCode);
  }

  http.end();
}

void setup() {
  Serial.begin(115200);
  delay(100);

  FastLED.addLeds<WS2812B, LED_PIN, GRB>(leds, NUM_LEDS);
  FastLED.setBrightness(BRIGHTNESS);
  FastLED.clear();
  FastLED.show();

  connectWiFi();
}

void loop() {
  fetchAndUpdate();
  delay(15000); // CTA data updates roughly every 15-20 seconds
}
\end{lstlisting}

\subsection{8.3 Configure the Sketch}

Edit the following lines in the sketch:

\begin{itemize}
    \item \texttt{ssid} and \texttt{password} to match your WiFi network.
    \item \texttt{ctaApiKey} to your CTA Train Tracker API key.
    \item \texttt{LED\_PIN} if you are using a different ESP32 GPIO pin. GPIO12 is a common choice, but avoid pins that conflict with boot modes.
    \item \texttt{BRIGHTNESS} to a value between 20 and 80 for low heat and low power.
\end{itemize}

\subsection{8.4 Upload and Test}

Select the correct ESP32 board and port in the Arduino IDE. Click Upload. Open the Serial Monitor at 115200 baud to see debug messages.

If everything works, the LEDs at the next stations of active trains should light in the appropriate line colors.

\section{Step 9: Finish the Frame}

\begin{enumerate}
    \item Add a diffuser layer between the LEDs and the printed map. Tracing paper or thin white acrylic works well.
    \item Mount the LED board inside a shadow box or deep frame.
    \item Use foam board spacers to create a gap between the map and the LEDs so the light spreads evenly.
    \item Bundle and secure excess wires on the back.
    \item Drill a hole for the power cable and route the 5V supply outside the frame.
    \item If the power supply is internal, ensure adequate ventilation and strain relief.
\end{enumerate}

\section{Step 10: Calibration and Troubleshooting}

\subsection{Common Issues}

\begin{longtable}{p{4.5cm} p{8.5cm}}
\toprule
\textbf{Symptom} & \textbf{Likely Cause / Fix} \\
\midrule
No LEDs light & Check power supply, common ground between ESP32 and LED strip, data direction. \\
First LED flickers or random colors & Use 74AHCT125 level shifter; shorten data wire; add 330~$\Omega$ resistor. \\
LEDs are dim or white-looking & Low power or voltage drop; add power injection; lower brightness. \\
API returns HTTP 403 & Invalid or inactive API key; check key and API access. \\
JSON parse error & Increase \texttt{DynamicJsonDocument} size; check URL string. \\
WiFi disconnects & Move ESP32 closer to router; use \texttt{WiFi.setAutoReconnect(true)}. \\
Some stations never light & Check that \texttt{stop\_id} in CSV matches CTA API \texttt{nextStaId}; check \texttt{led\_index}. \\
\bottomrule
\caption{Basic troubleshooting table.}
\label{tab:troubleshoot}
\end{longtable}

\subsection{Improving Shared Station Behavior}

Many Loop stations serve multiple lines. With one LED per station, only the last train written to the LED will be shown. Possible improvements:

\begin{itemize}
    \item Cycle colors for shared stations every 800~ms.
    \item Add a second LED at major shared stations.
    \item Blend colors using an occupancy array.
\end{itemize}

For a first build, the simple override behavior is acceptable.

\section{Safety and Legal Notes}

\begin{itemize}
    \item This project uses low voltage, but a 5V 10A supply can still melt wires if shorted. Use insulated connectors, avoid loose wires, and add a fuse if desired.
    \item Do not leave the device powered unattended until it has been tested for several hours.
    \item The CTA Train Tracker API is free for non-commercial use. Do not redistribute the data or use it for a commercial product without permission.
    \item CTA logos and official maps are protected. This guide is for personal/educational use.
\end{itemize}

\section{Possible Cost-Saving Substitutions}

\begin{itemize}
    \item Use a 5V 2A phone charger and set \texttt{BRIGHTNESS} to 40 or below. This will likely work if you never display full white.
    \item Use cardboard and black spray paint instead of foam board and a printed map.
    \item Skip the frame initially; hang the map with command strips.
    \item Use individual WS2812B LEDs purchased in bulk if you can find them cheaper than a full strip.
    \item Use a NodeMCU or Wemos D1 Mini if you already have one, but the ESP32 is recommended for HTTPS and RAM.
\end{itemize}

\section{Conclusion}

You should now have a functional, live-updating CTA wall map. The build is inexpensive, educational, and visually striking. Once the basic version works, you can add features such as:

\begin{itemize}
    \item A web portal for WiFi configuration.
    \item OTA firmware updates.
    \item A small OLED display showing train ETAs.
    \item Sound or vibration alerts.
    \item A custom PCB to replace point-to-point wiring.
\end{itemize}

Enjoy your build!

\appendix

\section{File Checklist}

\begin{table}[H]
\centering
\begin{tabular}{ll}
\toprule
\textbf{Filename} & \textbf{Purpose} \\
\midrule
\texttt{prepare\_gtfs.py} & Extract CTA rail stations from GTFS. \\
\texttt{stations\_template.csv} & Station list with empty LED index column. \\
\texttt{make\_mapping.py} & Generate \texttt{mapping.h} from filled CSV. \\
\texttt{mapping.h} & CTA \texttt{stop\_id} to LED index mapping. \\
\texttt{cta\_wall\_map.ino} & ESP32 Arduino firmware. \\
\bottomrule
\end{tabular}
\caption{Files used in this project.}
\label{tab:files}
\end{table}

\section{Complete Python Scripts}

The scripts are included in Listings~\ref{lst:prepare_gtfs} and \ref{lst:make_mapping}. Save them with the exact filenames and run them as described.

\section{Complete Arduino Sketch}

The Arduino sketch is included in Listing~\ref{lst:arduino}. Place it in a sketch folder with the generated \texttt{mapping.h}.

\end{document}